\documentclass[11pt]{article}

\usepackage[T1]{fontenc}
\usepackage{lmodern}
\usepackage[margin=30mm]{geometry}
\usepackage{microtype}
\usepackage{amsmath,amssymb,amsthm,mathtools}
\usepackage{xcolor}
\usepackage[colorlinks=true,linkcolor=blue!45!black,citecolor=blue!45!black,urlcolor=blue!55!black]{hyperref}

\newtheorem{theorem}{Theorem}
\newtheorem{proposition}[theorem]{Proposition}
\theoremstyle{definition}
\newtheorem{remark}[theorem]{Remark}
\newcommand{\rank}{\rho}
\newcommand{\layer}[1]{L_{#1}}

\title{A Counterexample for Rectangular Products\\
\large and a note on the asymmetric-channel literature}
\author{}
\date{August 2026}

\begin{document}
\maketitle

\begin{abstract}
We consider the direct rectangular analogue of the central-residue
construction for $k$-separated families in products of chains.  We give an
elementary counterexample for every $k\geq 2$, already when $n=2k+1$, the
first dimension not covered by the two-block argument for $k\geq n/2$.
We also record how the problem fits into Kova\v{c}evi\'c's framework for
asymmetric-error-detecting codes and note that this reference is not currently
present in the bibliography of the main paper.
\end{abstract}

\section{The natural rectangular construction}

Let
\[
  P(\mathbf b)=\prod_{i=1}^n\{0,1,\ldots,b_i\},
  \qquad \rank(x)=\sum_{i=1}^n x_i,
  \qquad T=\sum_{i=1}^n b_i.
\]
A family $\mathcal F\subseteq P(\mathbf b)$ is \emph{$k$-separated} if,
whenever $x<y$ coordinatewise are distinct members of $\mathcal F$, they
differ in more than $k$ coordinates.

Write $b_{(1)}\geq\cdots\geq b_{(n)}$ for the decreasing rearrangement of
the bounds, and put
\[
  D_k=\sum_{i=1}^k b_{(i)},\qquad M=D_k+1.
\]
The direct analogue of the two residue families in the equal-sided cube is
\[
  \mathcal A^-=
  \{x:\rank(x)\equiv\lfloor T/2\rfloor\pmod M\},\qquad
  \mathcal A^+=
  \{x:\rank(x)\equiv\lceil T/2\rceil\pmod M\}.
\]
Both are $k$-separated.  Indeed, if $x<y$ differ in at most $k$
coordinates, then
\[
  1\leq \rank(y)-\rank(x)\leq D_k<M.
\]
Coordinatewise reflection $x_i\mapsto b_i-x_i$ exchanges $\mathcal A^-$
and $\mathcal A^+$, so they have the same cardinality.

When $k\geq n/2$, one has $D_k\geq T/2$.  Thus $M>T/2$, and these residue
families reduce to the lower and upper central layers.  This is precisely the
range in which the two-block symmetric-chain argument extends to rectangular
products.  The construction below shows that the analogous statement fails
immediately outside that range.

\section{An infinite family of counterexamples}

For $r\in\mathbb Z$, write
\[
  \layer r=\{x\in P(\mathbf b):\rank(x)=r\}.
\]

\begin{theorem}
For every $k\geq2$, there is a rectangular product in dimension $2k+1$
for which a $k$-separated family is strictly larger than both central-residue
families $\mathcal A^-$ and $\mathcal A^+$.
\end{theorem}

\begin{proof}
Take $k+1$ coordinates of bound $k+1$ and $k$ coordinates of bound $k+2$:
\[
  P=\{0,\ldots,k+1\}^{k+1}\times
    \{0,\ldots,k+2\}^{k}.
\]
The sum of the $k$ largest bounds and the total rank are
\[
  D:=D_k=k(k+2),\qquad
  T=(k+1)^2+k(k+2)=2D+1.
\]
Consequently $M=D+1$, and the lower central-residue family is exactly
\[
  \mathcal A^-=\layer D\cup\{\mathbf b\},
\]
where $\mathbf b$ is the top vertex of $P$.

For each of the $k$ coordinates of bound $k+2$, let
$u_j=\mathbf b-e_j$ be the corresponding coatom, and define
\[
  \mathcal B=\layer D\cup\{u_1,\ldots,u_k\}.
\]
The layer $\layer D$ is an antichain, and the $u_j$ are pairwise
incomparable.  It remains only to rule out a forbidden comparison between
the two parts.

Suppose that $x\in\layer D$ satisfies $x\leq u_j$ and differs from $u_j$
in at most $k$ coordinates.  Since
$\rank(u_j)=2D$, the total decrease from $u_j$ to $x$ is $D$.  On the other
hand, the sum of the $k$ largest coordinates of $u_j$ is
\[
  (k-1)(k+2)+(k+1)=D-1.
\]
No change in at most $k$ coordinates can therefore decrease its rank by
$D$, a contradiction.  Hence $\mathcal B$ is $k$-separated.  Finally,
\[
  |\mathcal B|=|\layer D|+k
  >|\layer D|+1=|\mathcal A^-|=|\mathcal A^+|.
\]
\end{proof}

\begin{remark}[First instance]
For $k=2$, take bounds $(3,3,3,4,4)$.  Here $D=8$, $M=9$, and $T=17$.
The coefficient of $z^8$ in
\[
  (1+z+z^2+z^3)^3(1+z+z^2+z^3+z^4)^2
\]
is
\[
  \binom{12}{4}-3\binom84-2\binom74+3\binom44=218.
\]
Thus each central-residue family has size $219$, whereas
\[
  \layer8\cup\{(3,3,3,3,4),(3,3,3,4,3)\}
\]
is $2$-separated and has size $220$.  A maximum-independent-set computation
also gives optimum $220$, but that upper bound is not needed for the
counterexample.
\end{remark}

\begin{remark}[Uniqueness]
Even in the range $k\geq n/2$, the rectangular cardinality result does not
generally have the uniqueness conclusion of the equal-sided theorem.  For
$n=k=2$ and bounds $(1,2)$, the two central layers and
\[
  \{(1,0),(0,2)\}
\]
are three distinct maximum antichains of size $2$.
\end{remark}

\section{Relation with asymmetric-error-detecting codes}

Kova\v{c}evi\'c~\cite{kovacevic} studies codes in
$\{0,\ldots,q-1\}^n$ that detect coordinatewise-increasing errors subject
to three possible bounds: per-coordinate amplitude $a$, Hamming support
$h$, and total increase $t$.  The equal-sided problem in the main paper is
the finite, non-cyclic specialization
\[
  q=d+1,\qquad a=d,\qquad h=k,\qquad t=dk.
\]
Here the amplitude and total-increase bounds are redundant, so only the
support bound remains.  Thus Kova\v{c}evi\'c's \emph{model} contains the
problem, but his finite-alphabet optimality theorems do not contain the main
paper's result for $k>1$.  They establish the modular construction on the
infinite lattice and in several different finite parameter regimes, while
the unrestricted finite problem is left open.  The weighted-chain theorem
for $d=1,2$ therefore supplies an exact finite-boundary result not furnished
by that paper.

At the time of writing this note, Kova\v{c}evi\'c~\cite{kovacevic} is not
cited in the bibliography of the main paper.  It would be natural to add it
to the coding-theory discussion, while making clear that it provides the
ambient framework rather than the theorem proved here.

\begin{thebibliography}{1}
\bibitem{kovacevic}
M.~Kova\v{c}evi\'c,
\newblock On error detection in asymmetric channels,
\newblock \emph{IEEE Communications Letters} 21 (2017), no.~9, 1933--1936.
\newblock \url{https://arxiv.org/abs/1706.04540}.
\end{thebibliography}

\end{document}
